% Part: first-order-logic
% Chapter: tableaux
% Section: identity

\documentclass[../../../include/open-logic-section]{subfiles}

\begin{document}

\olfileid[hi]{fol}{tab}{ide}

\olsection{\usetoken{S}{identity} सहित \usetoken{P}{tableau}}

!!{identity} वाले !!^{tableau}s के लिए अतिरिक्त अनुमान-नियम चाहिए। $\eq$ के
नियम ये हैं (यहाँ $t$, $t_1$ और $t_2$ बंद पद हैं):

\begin{defish}
\AxiomC{}
\RightLabel{$\eq$}
\UnaryInfC{\sFmla{\True}{\eq[t][t]}}
\DisplayProof
\hfill
\AxiomC{\sFmla{\True}{\eq[t_1][t_2]}}
\noLine
\UnaryInfC{\sFmla{\True}{!A(t_1)}}
\RightLabel{$\TRule{\True}{\eq}$}
\UnaryInfC{\sFmla{\True}{!A(t_2)}}
\DisplayProof
\hfill
\AxiomC{\sFmla{\True}{\eq[t_1][t_2]}}
\noLine
\UnaryInfC{\sFmla{\False}{!A(t_1)}}
\RightLabel{$\TRule{\False}{\eq}$}
\UnaryInfC{\sFmla{\False}{!A(t_2)}}
\DisplayProof
\end{defish}
ध्यान दें कि अन्य सभी नियमों के विपरीत, $\TRule{\True}{\eq}$ और
$\TRule{\False}{\eq}$ के लिए शाखा पर \emph{दो} चिह्नित !!{formula}s पहले
से उपस्थित होने चाहिए, अर्थात् $\sFmla{\True}{\eq[t_1][t_2]}$ और
$\sFmla{S}{!A(t_1)}$ दोनों।

\begin{ex}
यदि $s$ और $t$ बंद पद हैं, तो $\eq[s][t], !A(s)
\Proves !A(t)$:
\begin{oltableau}
  [\sFmla{\False}{\formula{A}(t)}, just = \TAss
    [\sFmla{\True}{\eq[s][t]}, just = \TAss
      [\sFmla{\True}{\formula{A}(s)}, just = \TAss
        [\sFmla{\True}{\formula{A}(t)}, just={\TRule{\True}{\eq}[2, 3]}, close]
      ]
    ]
  ]
\end{oltableau}
इसे समान वस्तुओं की प्रतिस्थापनीयता का सिद्धांत, अथवा लाइबनित्स का नियम,
कहा जाता है।

!!^{tableau}s सिद्ध करते हैं कि $\eq$ सममित है, अर्थात्
$\eq[s_1][s_2]
\Proves \eq[s_2][s_1]$:
\begin{oltableau}
  [\sFmla{\False}{\eq[s_2][s_1]}, just = \TAss
    [\sFmla{\True}{\eq[s_1][s_2]}, just = \TAss
      [\sFmla{\True}{\eq[s_1][s_1]}, just = {$\eq$}
        [\sFmla{\True}{\eq[s_2][s_1]}, just = {\TRule{\True}{\eq}[2, 3]}, close]
      ]
    ]
  ]
\end{oltableau}
यहाँ पंक्ति~$2$ पहला अपेक्षित
!!{formula}~$\sFmla{\True}{\eq[s_1][s_2]}$ है, जिसकी आवश्यकता
$\TRule{\True}{\eq}$ को है। पंक्ति~$3$ दूसरा है, जो
$\sFmla{\True}{!A(s_2)}$ के रूप का है---$!A(x)$ को $\eq[x][s_1]$ समझें;
तब $!A(s_1)$, $\eq[s_1][s_1]$ और $!A(s_2)$, $\eq[s_2][s_1]$ है।

वे यह भी सिद्ध करते हैं कि $\eq$ संक्रामक है, अर्थात्
$\eq[s_1][s_2],
\eq[s_2][s_3] \Proves \eq[s_1][s_3]$:
\begin{oltableau}
  [\sFmla{\False}{\eq[s_1][s_3]}, just = \TAss
    [\sFmla{\True}{\eq[s_1][s_2]}, just = \TAss
      [\sFmla{\True}{\eq[s_2][s_3]}, just = \TAss
        [\sFmla{\True}{\eq[s_1][s_3]}, just = {\TRule{\True}{\eq}[3, 2]}, close]
      ]
    ]
  ]
\end{oltableau}
इस !!{tableau} में $\TRule{\True}{\eq}$ का पहला अपेक्षित !!{formula}
पंक्ति~$3$, $\sFmla{\True}{\eq[s_2][s_3]}$ है ($s_2$,~$t_1$ की भूमिका
और $s_3$,~$t_2$ की भूमिका निभाता है)। दूसरा अपेक्षित सूत्र, जो
$\sFmla{\True}{!A(s_2)}$ के रूप का है, पंक्ति~$2$ है। यहाँ $!A(x)$ को
$\eq[s_1][x]$ समझें; तब $!A(s_2)$, $\eq[t_1][t_2]$ (अर्थात् पंक्ति~$2$)
और $!A(s_3)$, निष्कर्ष का !!{formula}~$\eq[s_1][s_3]$ बनता है।
\end{ex}

\begin{prob}
निम्न के लिए बंद !!{tableau}s दीजिए:
\begin{enumerate}
\item $\sFmla{\False}{\lforall[x][\lforall[y][((x = y \land !A(x))
      \lif !A(y))]]}$
\item $\sFmla{\False}{\lexists[x][(!A(x) \land
    \lforall[y][(!A(y) \lif y = x)])]}$,\\
  $\sFmla{\True}{\lexists[x][!A(x)] \land
    \lforall[y][\lforall[z][((!A(y) \land !A(z)) \lif y = z)]]}$
\end{enumerate}
\end{prob}

\end{document}
